\documentclass[11pt,letterpaper]{article}

% ── geometry & layout ────────────────────────────────────────────────
\usepackage[margin=1in]{geometry}
\usepackage{fancyhdr}
\usepackage{setspace}
\onehalfspacing

% ── math ─────────────────────────────────────────────────────────────
\usepackage{amsmath,amssymb,amsthm,mathtools}
\usepackage[T1]{fontenc}

% ── colour & boxes ───────────────────────────────────────────────────
\usepackage[dvipsnames]{xcolor}
\usepackage[most]{tcolorbox}
\usepackage{enumitem}
\usepackage{booktabs,array,tabularx}
\usepackage{tikz}
\usetikzlibrary{arrows.meta,positioning,calc,decorations.pathreplacing,matrix}
\usepackage[round,authoryear]{natbib}
\usepackage[colorlinks=true,linkcolor=MidnightBlue,citecolor=OliveGreen,urlcolor=BrickRed]{hyperref}

% ── theorem environments ─────────────────────────────────────────────
\theoremstyle{definition}
\newtheorem{definition}{Definition}[section]
\newtheorem{example}{Example}[section]
\theoremstyle{plain}
\newtheorem{theorem}{Theorem}[section]
\newtheorem{proposition}{Proposition}[section]
\newtheorem{lemma}{Lemma}[section]
\newtheorem{corollary}{Corollary}[section]
\theoremstyle{remark}
\newtheorem*{remark}{Remark}

% ── macros: PETERS-CANONICAL (see crosswalk, App. A) ──────────────────
\newcommand{\R}{\mathbb{R}}
\newcommand{\Q}{\mathbb{Q}}
\newcommand{\E}{\mathbb{E}}
\newcommand{\voters}{N}            % set of voters
\newcommand{\projects}{C}          % set of projects (candidates)
\newcommand{\budget}{b}            % total budget
\newcommand{\votes}{v}             % per-voter vote budget (DEM ballot width)
\newcommand{\cost}{\mathrm{cost}}  % cost function, written cost(c)
\newcommand{\Dist}{\mathcal{D}}    % set of districts
\newcommand{\salience}{s}          % district salience  s_d   (was \alpha_d)
\newcommand{\info}{\kappa}         % voter information fn  kappa_i  (was \sigma_i)
\newcommand{\popshare}{\pi}        % district population share  pi_d  (was \beta_d)
\newcommand{\DEM}{\mathrm{DEM}}
\newcommand{\UAM}{\mathrm{UAM}}
\newcommand{\MES}{\mathrm{MES}}
\newcommand{\EJR}{\mathrm{EJR}}
\newcommand{\FJR}{\mathrm{FJR}}
\newcommand{\PJR}{\mathrm{PJR}}
\newcommand{\JRax}{\mathrm{JR}}
\newcommand{\GCR}{\mathrm{GCR}}
\newcommand{\PAV}{\mathrm{PAV}}
\newcommand{\PSC}{\mathrm{PSC}}
\DeclareMathOperator*{\argmin}{arg\,min}
\DeclareMathOperator*{\argmax}{arg\,max}

% ── pedagogical boxes (identical to notebook_v2) ──────────────────────
\newtcolorbox{keyidea}{colback=OliveGreen!6,colframe=OliveGreen!65,
  fonttitle=\bfseries,title={Key idea},boxsep=2pt,left=5pt,right=5pt,top=3pt,bottom=3pt}
\newtcolorbox{intuition}{colback=MidnightBlue!5,colframe=MidnightBlue!60,
  fonttitle=\bfseries,title={Intuition},boxsep=2pt,left=5pt,right=5pt,top=3pt,bottom=3pt}
\newtcolorbox{connection}{colback=TealBlue!6,colframe=TealBlue!60,
  fonttitle=\bfseries,title={Connection to Durango},boxsep=2pt,left=5pt,right=5pt,top=3pt,bottom=3pt}
\newtcolorbox{naming}{colback=Violet!6,colframe=Violet!65,
  fonttitle=\bfseries,title={A name we are introducing (our label, not standard terminology)},
  boxsep=2pt,left=5pt,right=5pt,top=3pt,bottom=3pt}
\newtcolorbox{watchout}{colback=BrickRed!5,colframe=BrickRed!60,
  fonttitle=\bfseries,title={Watch out --- a subtle point that is easy to get wrong},
  boxsep=2pt,left=5pt,right=5pt,top=3pt,bottom=3pt}
\newtcolorbox{aside}{colback=Dandelion!8,colframe=Dandelion!80!black,
  boxsep=2pt,left=5pt,right=5pt,top=3pt,bottom=3pt}

% ── header ───────────────────────────────────────────────────────────
\pagestyle{fancy}\fancyhf{}
\fancyhead[L]{\small\textit{Foundations Notebook --- built on Peters, Pierczy\'nski \& Skowron}}
\fancyhead[R]{\small\textit{Equal Shares for Durango PB}}
\fancyfoot[C]{\thepage}
\renewcommand{\headrulewidth}{0.4pt}

\title{%
  \vspace{-0.8cm}
  {\large\scshape Foundations Notebook --- canonical base}\\[0.6em]
  {\LARGE\bfseries The Method of Equal Shares,\\[0.2em] founded on Peters--Pierczy\'nski--Skowron,\\[0.2em] applied to Durango}\\[0.5em]
  {\large Notation, definitions, and theory re-derived on the rock of\\[0.2em]
  \emph{Proportional Participatory Budgeting with Additive Utilities}}
}
\author{\textbf{Mario Alberto Garc\'ia Meza}\\[2pt]
  \small Facultad de Econom\'ia, Contadur\'ia y Administraci\'on\\
  \small Universidad Ju\'arez del Estado de Durango (FECA--UJED)}
\date{June 2026 \;\textbullet\; Tight first cut (Sections 1--7 complete; Section 8 a stub)}

\begin{document}
\maketitle
\thispagestyle{fancy}

\begin{abstract}
\noindent This notebook re-founds our entire participatory-budgeting (PB) theory on a single
source: Peters, Pierczy\'nski \& Skowron's \emph{Proportional Participatory Budgeting with
Additive Utilities} (arXiv:2008.13276v2; NeurIPS~2021). We adopt their notation and definitions
wholesale, reproduce the load-bearing results (the Method of Equal Shares, Extended and Fully
Justified Representation, the core, priceability), and then rebuild our Durango contribution on
top: the District-Exclusivity Mechanism (DEM), its distortions, and the case that Equal Shares
repairs them while delivering geographic equity \emph{by construction}. Remarkably, Peters et
al.'s own motivating example---\emph{Circleville}, a city run as separate district elections---is
our Durango critique almost verbatim. This is the canonical reference for the project; its section
order is the intended skeleton of the paper.
\end{abstract}

\begin{aside}
\small\textbf{How to read this, and three conventions.}
\textbf{(1) Boxes.} Green = the load-bearing idea; blue = intuition; teal = the bridge to Durango;
\emph{red ``Watch out''} = a subtle point that is genuinely easy to get wrong; violet = a label we
are coining (so you can always tell ours from the literature's).
\textbf{(2) Provenance.} The Method of Equal Shares, EJR, FJR, the core, priceability, PAV, PSC,
and GCR are \emph{theirs} and are cited. The District-Exclusivity Mechanism, the Unrestricted
Allocation Mechanism, district \emph{salience}, the voter \emph{information function}, and
\emph{$\popshare_d$-geographic equity} are \emph{ours} and are flagged in violet boxes.
\textbf{(3) Notation.} We now follow Peters et al.\ exactly. Appendix~\ref{app:crosswalk} gives the
crosswalk from our old \texttt{notebook\_v2} symbols; the headline changes are
$\projects$ for projects, $\budget$ for budget, $\cost(\cdot)$ for cost, plain $\rho$ for the
Equal-Shares threshold, and the renamings $s_d$ (salience), $\kappa_i$ (information), $\pi_d$
(population share) that free $\alpha,\sigma,\beta$ for their roles in Peters et al.
\end{aside}

\tableofcontents
\newpage

% ═════════════════════════════════════════════════════════════════════
\section{Introduction: from Circleville to Durango}
\label{sec:intro}

A growing number of cities allocate part of their budget by \emph{participatory budgeting} (PB):
residents vote over candidate projects, and a rule turns the ballots into a funded set subject to
a budget constraint (\citealp{cabannes2004}; \citealp{wampler2021}; \citealp{azizshah2020}). The
question this project asks is narrow and practical: \emph{is the rule Durango uses a good one, and
if not, what should replace it?}

Peters, Pierczy\'nski \& Skowron (2021) open with a fictional city, \textbf{Circleville}, split into
four districts of unequal size. If projects are local and residents vote only for their own
district's projects, then a single city-wide ``most-votes-win'' count hands the entire budget to
the largest district---what they call \emph{the tyranny of the largest minority}. The standard fix
is to \emph{run separate elections per district}, dividing the budget in advance (often in
proportion to population) and letting residents vote only locally. Peters et al.\ then catalogue
exactly what that fix breaks:
\begin{itemize}[leftmargin=1.4em,itemsep=2pt,topsep=2pt]
  \item \textbf{Boundary projects are orphaned.} A project on the border of two districts must be
  assigned to one; residents of the other cannot vote for it, so it is under-funded relative to its
  true value.
  \item \textbf{City-wide goods are hard to place.} A project benefiting the whole city has no
  natural district.
  \item \textbf{Non-geographic interest groups are silenced.} Parents, or cyclists, spread across
  the city cannot form a voting bloc, even when their shared project would be highly valuable.
\end{itemize}

\begin{connection}
Durango's \emph{presupuesto participativo} \textbf{is Circleville's ``separate district
elections.''} A resident selects one district and casts a fixed number of votes
($\votes$, e.g.\ six) \emph{within} that district only. This is precisely the fix Peters et al.\
diagnose---so their critique is, line for line, our motivation. Our contribution is to (i) make the
distortions precise as propositions in their model, with Durango's institutional detail, and (ii)
argue that their remedy, the Method of Equal Shares, is the right reform for Durango because it
keeps a single city-wide election while protecting marginalized districts \emph{by construction}.
\end{connection}

Two features of the Durango setting are not incidental; they shape the whole problem and must stay
central.

\begin{naming}
\textbf{The voter information function $\kappa_i$ (rational ignorance).} A voter typically knows
only the projects in her immediate circle. Under Durango's rule she must nonetheless spend
\emph{all} $\votes$ votes inside her chosen district---including on projects she has never heard of.
We write $\kappa_i \subseteq \projects$ for the set of projects voter $i$ can actually evaluate.
The mechanism forces votes outside $\kappa_i$, converting ignorance into noise. (Background:
rational ignorance, Downs 1957; costly information in voting, Martinelli 2006.)
\end{naming}

\begin{naming}
\textbf{District salience $s_d$ and the centralization risk.} Removing the district constraint
entirely (our \emph{Unrestricted Allocation Mechanism}, \S\ref{sec:model}) lets high-visibility
city-center projects absorb the vote. We capture visibility by a \emph{salience} parameter $s_d$
per district: high-$s_d$ districts attract attention, low-$s_d$ (often marginalized) districts get
none. The district constraint is a \emph{blunt equity instrument}---it protects low-salience
districts by accident. The reform must preserve that protection deliberately, not discard it.
\end{naming}

\begin{keyidea}
The design target Peters et al.\ identify, and which we adopt for Durango, is: \emph{a single
city-wide election whose rule spends money proportionally}, automatically identifying groups with
common interests (geographic or not) and guaranteeing each its fair share. That target is made
precise by the axiom \emph{Extended Justified Representation} (EJR) and met by the \emph{Method of
Equal Shares}.
\end{keyidea}

The rest of the notebook follows Peters et al.'s architecture: the model
(\S\ref{sec:model}), the Method of Equal Shares (\S\ref{sec:mes}), the proportionality axioms
(\S\ref{sec:axioms}), then our Durango layer---the DEM's distortions (\S\ref{sec:distortions}),
their repair under Equal Shares (\S\ref{sec:repair}), the three-way comparison
(\S\ref{sec:comparison}), and the simulation design (\S\ref{sec:simulation}).

% ═════════════════════════════════════════════════════════════════════
\section{The participatory-budgeting model}
\label{sec:model}

We use Peters et al.'s model verbatim, then add the few district-level objects Durango needs.
For $t\in\mathbb{N}$ write $[t]=\{1,\dots,t\}$.

\begin{definition}[Election; general PB model]
An \emph{election} is a tuple $E=(\voters, \projects, \budget, \cost, \{u_i\}_{i\in \voters})$ where
\begin{itemize}[leftmargin=1.4em,itemsep=1pt,topsep=2pt]
  \item $\voters=[n]$ is the set of \emph{voters} and $\projects=\{c_1,\dots,c_m\}$ the set of
  \emph{projects} (candidates);
  \item $\budget\in\Q_{\ge 0}$ is the available \emph{budget};
  \item $\cost\colon \projects\to\Q_{\ge 0}$ gives each project's cost, extended to sets by
  $\cost(T)=\sum_{c\in T}\cost(c)$;
  \item each $u_i\colon \projects\to\Q_{\ge 0}$ is voter $i$'s \emph{additive} utility, so a funded
  set $T$ gives $i$ utility $u_i(T)=\sum_{c\in T}u_i(c)$; for $S\subseteq \voters$ write
  $u_S(T)=\sum_{i\in S}u_i(T)$. We assume $u_\voters(c)>0$ for every $c$.
\end{itemize}
A set $W\subseteq \projects$ is \emph{feasible} if $\cost(W)\le \budget$; a feasible $W$ is an
\emph{outcome}. An \emph{(aggregation) rule} $R$ maps each election $E$ to an outcome $R(E)$.
\end{definition}

\begin{definition}[Two special cases]
\emph{Committee elections / unit costs:} $\budget\in\mathbb{N}$ and $\cost(c)=1$ for all $c$, so
$W$ is feasible iff $|W|\le \budget$ (a committee of size $\budget$).
\emph{Approval utilities:} $u_i(c)\in\{0,1\}$; then $A(i)=\{c: u_i(c)=1\}$ is $i$'s
\emph{approval set}, and $c\in A(i)\cap W$ is a \emph{representative} of $i$ in $W$.
\end{definition}

\begin{watchout}
Two utility readings of the \emph{same} approval ballot differ once costs vary. Under
\emph{approval utilities} a voter's satisfaction is the \emph{number} of her approved projects that
get funded ($u_i(c)=1$). Under \emph{cost utilities} it is the total \emph{money} spent on them
($u_i(c)=\cost(c)$). Peters et al.\ note the latter is often closer to the truth (their Paris
example: a project $186\times$ more expensive but only $1.15\times$ more popular). Which reading we
use is a modeling choice we must make explicit in the simulation (\S\ref{sec:simulation}).
\end{watchout}

\subsection{Durango-specific objects}

To talk about Durango we layer a district structure on top of the general model.

\begin{definition}[District structure]
Let $\Dist$ be a set of \emph{districts}. A map $\delta\colon \projects\to\Dist$ assigns each
project to a district, and each voter $i$ resides in a district $d(i)\in\Dist$. Write
$n_d=|\{i: d(i)=d\}|$ for the population of district $d$, and $\popshare_d=n_d/n$ for its
\emph{population share}. A pre-divided budget assigns $\budget_d$ to district $d$ with
$\sum_d \budget_d=\budget$ (Durango uses $\budget_d \approx \popshare_d\,\budget$).
\end{definition}

\begin{naming}
\textbf{District salience} $s_d\in\R_{>0}$ scales how much attention/turnout district $d$'s projects
attract in an \emph{unconstrained} election. \textbf{The information function} $\kappa_i\subseteq
\projects$ is the set of projects voter $i$ can evaluate. Both are ours, not Peters et al.'s.
\end{naming}

\begin{naming}
\textbf{The District-Exclusivity Mechanism ($\DEM$).} Each voter $i$ selects exactly one district
$d_i\in\Dist$ and allocates her $\votes$ votes among projects in $d_i$ \emph{only}. Within each
district, projects are funded in decreasing order of total votes until $\budget_d$ is exhausted.
This is Durango's actual rule and Peters et al.'s ``separate district elections.''
\end{naming}

\begin{naming}
\textbf{The Unrestricted Allocation Mechanism ($\UAM$).} The same ballot but with no district
restriction: a voter may direct her $\votes$ votes at \emph{any} projects, and projects are funded
in decreasing order of total votes against the single global budget $\budget$. This is the naive
``most-votes-win'' city-wide rule---the one that yields the tyranny of the largest (most salient)
minority.
\end{naming}

\begin{intuition}
$\DEM$ and $\UAM$ are the two poles Peters et al.\ warn about: $\UAM$ centralizes (high-$s_d$
districts win everything), $\DEM$ fragments (the per-district fix that orphans boundary projects,
silences city-wide groups, and---via $\kappa_i$---forces uninformed votes). The Method of Equal
Shares is the third option that avoids both poles.
\end{intuition}

% ═════════════════════════════════════════════════════════════════════
\section{The Method of Equal Shares}
\label{sec:mes}

\begin{naming}
The \emph{Method of Equal Shares} ($\MES$) is \textbf{not} ours; it is Peters et al.'s rule
(originally ``Rule~X,'' Peters \& Skowron 2020). We restate it exactly.
\end{naming}

\begin{definition}[Method of Equal Shares]
\label{def:mes}
Give each voter an equal share $\budget/n$ of the budget. Start with $W=\varnothing$ and add
projects one at a time. For a chosen $c$, let $p_i(c)\ge 0$ be $i$'s payment, with
$\sum_{i\in \voters}p_i(c)=\cost(c)$; write $p_i(W)=\sum_{c\in W}p_i(c)$ and the remaining balance
$b_i=\budget/n-p_i(W)$. For $\rho\ge 0$, a project $c\notin W$ is \emph{$\rho$-affordable} if
\[
  \sum_{i\in \voters}\min\bigl(b_i,\; u_i(c)\cdot \rho\bigr)=\cost(c).
\]
If no project is $\rho$-affordable for any $\rho$, $\MES$ terminates and returns $W$. Otherwise it
selects a $c\notin W$ that is $\rho$-affordable for the \emph{minimum} $\rho$, with payments
$p_i(c)=\min(b_i,\,u_i(c)\cdot\rho)$.
\end{definition}

\begin{keyidea}
Each voter pays \emph{per unit of her utility}, at a common price $\rho$, capped at her remaining
balance. Projects enter in increasing order of $\rho$, i.e.\ in \emph{decreasing order of utility
per dollar}. So $\MES$ favours cheap projects, projects with many supporters, and projects whose
supporters still hold money (those whose preferences are so far least satisfied).
\end{keyidea}

\begin{watchout}
We now write the threshold as plain $\rho$, matching Peters et al. In the old \texttt{notebook\_v2}
it was $\varrho$ (\texttt{\textbackslash mesrho}) to dodge a LaTeX clash; that dodge is retired.
\end{watchout}

\paragraph{Approval special case.} With $u_i(c)\in\{0,1\}$, $c$ is $\rho$-affordable iff its cost
can be covered by its approvers so that each pays at most $\rho$: approvers with less than $\rho$
left spend all they have, the rest pay exactly $\rho$. Cost is split \emph{as equally as
possible}---hence ``Equal Shares.'' The first project chosen is the affordable one maximizing
\emph{approvers per dollar}.

\paragraph{Computing $\rho$ (Algorithm 1).} For $c\notin W$ with supporters $i_1,\dots,i_t$ (those
with $u_i(c)>0$), sort by $b_{i}/u_{i}(c)$ ascending. Supporters who exhaust their balance attain
the cap in the first coordinate. Iterating $s=1,\dots,t$, the candidate value is
\[
  \rho \;=\; \frac{\cost(c)-\bigl(b_{i_1}+\cdots+b_{i_{s-1}}\bigr)}{u_{i_s}(c)+\cdots+u_{i_t}(c)},
\]
accepted as soon as $\rho\cdot u_{i_s}(c)\le b_{i_s}$. The rule then picks
$c=\argmin_{c\notin W}\rho(c)$. With running partial sums this is $O(m^2 n\log n)$.

\begin{example}[Equal Shares, worked (Peters et al.'s instance)]
Ten voters, budget $\budget=\$100$ (so each starts with $\$10$), approval utilities, five projects:
P1 $\$36$ approved by $i_1$--$i_5$; P2 $\$36$ by $i_1$--$i_6$; P3 $\$25$ by $i_1$--$i_5$;
P4 $\$24$ by $i_4$--$i_7$; P5 $\$24$ by $i_6$--$i_9$. Round 1 thresholds: $\rho_{P1}=7.2$,
$\rho_{P2}=6$, $\rho_{P3}=5$, $\rho_{P4}=\rho_{P5}=6$. $\MES$ funds \textbf{P3} (lowest $\rho$);
$i_1$--$i_5$ each pay $\$5$. Now P1, P2 are unaffordable (supporters too poor); recomputing,
$\rho_{P4}=7$, $\rho_{P5}=6$, so $\MES$ funds \textbf{P5}; $i_6$--$i_9$ pay $\$6$. Nothing else is
affordable: $W=\{\text{P3},\text{P5}\}$, spending only $\$49$ of $\$100$. The leftover illustrates
\emph{non-exhaustiveness} (\S\ref{sec:axioms}); completion methods address it.
\end{example}

\begin{connection}
Read the five projects as Durango projects from different districts and the ten voters as a mixed
electorate. $\MES$ never asked anyone to ``spend all six votes in one district.'' A voter pays only
for projects she actually values---exactly what $\kappa_i$ demands---and a small, cohesive group
(e.g.\ supporters of P5) secures its project out of \emph{its own} share, with no district quota
imposed from above.
\end{connection}

% ═════════════════════════════════════════════════════════════════════
\section{Proportionality: EJR, the core, priceability, FJR}
\label{sec:axioms}

\subsection{Why not just maximize welfare? (Theorem 1)}

\begin{theorem}[Costs matter; \citealp{peters2021}]
\label{thm:welfarism}
Every rule that depends only on voters' utility functions and on the collection of budget-feasible
sets must fail proportionality---even on instances with a clean district structure.
\end{theorem}

\begin{intuition}
\emph{Onetown vs.\ Twotown.} Two cities with identical voters, districts, and approval sets, and
even identical feasibility constraints, but different \emph{costs}. Proportional Approval Voting
(PAV), which maximizes $\sum_i H(|A(i)\cap W|)$ with $H$ the harmonic number, must return the same
outcome in both---yet the proportional outcome differs between them. Hence any welfare-optimizing
rule that ignores costs cannot be proportional once costs vary. This is why we cannot simply ``run
PAV'' on Durango, and why a \emph{cost-aware, sequential} rule like Equal Shares is needed.
\end{intuition}

\subsection{Extended Justified Representation (EJR)}

EJR formalizes ``no cohesive group is under-served.'' Peters et al.\ build it in three steps.

\begin{definition}[EJR, committee/unit-cost; \citealp{aziz2017}]
A group $S$ is \emph{$\ell$-cohesive} if $|S|\ge \tfrac{\ell}{\budget}\,n$ and
$\bigl|\bigcap_{i\in S}A(i)\bigr|\ge \ell$. A rule satisfies \emph{EJR} if for every $\ell$ and
every $\ell$-cohesive $S$ there is $i\in S$ with $|A(i)\cap R(E)|\ge \ell$.
\end{definition}

\begin{definition}[EJR, approval, general costs]
$S$ is \emph{$T$-cohesive} ($T\subseteq \projects$) if $|S|/n\ge \cost(T)/\budget$ and
$T\subseteq\bigcap_{i\in S}A(i)$. EJR requires some $i\in S$ with $|A(i)\cap R(E)|\ge |T|$.
\end{definition}

\begin{definition}[EJR, general additive]
\label{def:ejr}
$S$ is \emph{$(\alpha,T)$-cohesive}, for $\alpha\colon \projects\to[0,1]$ and $T\subseteq
\projects$, if $|S|/n\ge \cost(T)/\budget$ and $u_i(c)\ge \alpha(c)$ for all $i\in S,\,c\in T$. A
rule satisfies \emph{EJR} if for every such $S$ there is $i\in S$ with
$u_i(R(E))\ge \sum_{c\in T}\alpha(c)$.
\end{definition}

\begin{keyidea}
A group that is an $|S|/n$ fraction of voters and can name an affordable bundle $T$ they all value
is ``owed'' the utility of $T$---and EJR guarantees \emph{at least one} member actually gets it.
The ``at least one'' is not a weakness: requiring it for \emph{every} member is impossible even on
tiny instances (\citealp{aziz2017}).
\end{keyidea}

\begin{watchout}
EJR is genuinely demanding here. With a single voter, an EJR outcome must solve the
\emph{knapsack} problem; since knapsack is (weakly) NP-hard, \emph{no strongly polynomial rule can
satisfy EJR exactly in the general additive model} (Peters et al., Prop.~1). Equal Shares is
polynomial, so it cannot satisfy EJR exactly there---only ``up to one project.''
\end{watchout}

\begin{definition}[EJR up to one project]
\label{def:ejr1}
$R$ satisfies \emph{EJR up to one project} if for every $(\alpha,T)$-cohesive $S$ there is
$i^\ast\in S$ such that either $u_{i^\ast}(R(E))\ge\sum_{c\in T}\alpha(c)$, or for some $c^\ast\in
T$, $u_{i^\ast}(R(E)\cup\{c^\ast\})>\sum_{c\in T}\alpha(c)$.
\end{definition}

\begin{theorem}[\citealp{peters2021}]
\label{thm:mes-ejr}
The Method of Equal Shares satisfies EJR up to one project in the general PB model. In the
approval-based model (any costs) the ``up to one'' clause never bites, so $\MES$ satisfies EJR
\emph{exactly}.
\end{theorem}

\begin{proof}[Proof idea (the three-runs argument)]
Fix an $(\alpha,T)$-cohesive $S$; set $\alpha_c=\min_{i\in S}u_i(c)$ and $\alpha=\sum_{c\in
T}\alpha_c$. Compare three runs of $\MES$: (A) the real run, output $W$; (B) the run in which
voters of $S$ face \emph{no} budget cap when buying projects of $T$; (C) a shrunken run with only
$S$ and only $T$, unlimited budgets, and $u_i(c)=\alpha_c$. In (C) each $c\in T$ is
$\sigma_c$-affordable with $\sigma_c=\cost(c)/(|S|\alpha_c)$. Track, for a fixed voter $i^\ast$, the
money $f_{(B)}(x)$ and $f_{(C)}(x)$ needed to reach utility $x$. One shows
$f_{(C)}(x)\ge f_{(B)}(x)$ (the slope of $f_{(C)}$, namely $\sigma_{c_s}$, dominates that of
$f_{(B)}$, because in (B) there are extra payers and weakly higher utilities). Since
$f_{(C)}(\alpha)=\cost(T)/|S|\le \budget/n$, we get $f_{(B)}(\alpha)\le \budget/n$: by the time
$i^\ast$ has spent her whole share she has reached utility $\alpha$ (giving the exact bound), or
overshoots on one extra project $c^\ast\in T$ (the ``up to one'' case). \hfill
\end{proof}

\begin{corollary}[District proportionality, the easy case]
\label{cor:district}
If every voter in district $d$ approves exactly the projects of $d$, they form a $T$-cohesive group
for $T$ any affordable bundle of $d$'s projects with $\cost(T)\le \popshare_d\,\budget$. By
Theorem~\ref{thm:mes-ejr}, $\MES$ funds at least $|T|$ of them for some member---so district $d$
controls at least its population share $\popshare_d$ of the budget.
\end{corollary}

\begin{connection}
Corollary~\ref{cor:district} is the heart of the Durango argument. The DEM \emph{imposes}
$\budget_d\approx\popshare_d\,\budget$ from above (a quota). $\MES$ \emph{produces} the same
$\popshare_d$ floor \emph{endogenously}, while \emph{also} representing non-geographic groups
(parents, cyclists) and never forcing an uninformed vote. We name this guarantee
$\popshare_d$-\emph{geographic equity} and develop it in \S\ref{sec:repair}.
\end{connection}

\subsection{Stronger and weaker neighbours: the JR lattice}

\begin{definition}[The core]
$W$ is in the \emph{core} if for every $S\subseteq \voters$ and $T\subseteq \projects$ with
$|S|/n\ge\cost(T)/\budget$ there is $i\in S$ with $u_i(W)\ge u_i(T)$.
\end{definition}

The core lets a group propose \emph{any} affordable counter-bundle (not just a unanimously liked
one), so it implies EJR. It may be \emph{empty}, even with unit costs (Peters et al., Example~2),
which is why EJR---restricting to cohesive $T$---is the workable target. Equal Shares nonetheless
\emph{approximates} the core:

\begin{theorem}[\citealp{peters2021}]
With $u_{\max}/u_{\min}$ the ratio of the largest to smallest positive attainable voter utilities,
the Equal-Shares outcome lies in the $\alpha$-core for $\alpha=4\log(2\,u_{\max}/u_{\min})$.
\end{theorem}

\begin{definition}[Priceability; \citealp{petersskowron2020}]
\label{def:price}
$W$ is \emph{priceable} if there is a budget $b'\ge\budget$ (each voter holding $b'/n$) and
payments $p_i\colon \projects\to\R_{\ge0}$ with: (C1) $u_i(c)=0\Rightarrow p_i(c)=0$;
(C2) $\sum_c p_i(c)\le b'/n$; (C3) each $c\in W$ is fully paid; (C4) nothing is paid for $c\notin
W$; (C5) for each $c\notin W$ the unspent money of its supporters is $\le\cost(c)$. Equal Shares is
always priceable---it builds the price system as it runs.
\end{definition}

\begin{definition}[Exhaustiveness; \citealp{aziztalmon2018}]
$R$ is \emph{exhaustive} if no unfunded project fits in the remaining budget:
$\cost(R(E)\cup\{c\})>\budget$ for every $c\notin R(E)$.
\end{definition}

\begin{watchout}
Equal Shares is \emph{not} exhaustive: a project can remain affordable in the aggregate while its
own supporters are individually too poor to buy it, and the rule then declines to fund it.
Moreover exhaustiveness is \emph{incompatible} with priceability in general (Peters et al.,
Example~3). Completion methods (next) restore spending when that is desired.
\end{watchout}

\paragraph{Completion methods (for the leftover budget).}
\begin{itemize}[leftmargin=1.4em,itemsep=2pt,topsep=2pt]
  \item \emph{Varying the budget:} rerun $\MES$ with a gradually increased virtual budget $b'\ge
  \budget$, stopping just before real spending would exceed $\budget$. Keeps equal voting power and
  priceability; mirrors the d'Hondt divisor method.
  \item \emph{Perturbation:} set tiny $u_i(c)=\varepsilon$ wherever utility was $0$; with all
  utilities positive, $\MES$ is exhaustive (Peters et al., Prop.~2). Variant: $\varepsilon\cdot
  \cost(c)$, which experiments suggest is preferable.
  \item \emph{Utilitarian greedy:} spend the remainder on unfunded projects by total utility
  $\sum_i u_i(c)$---this mimics the plurality rule cities already use.
\end{itemize}

\subsection{Fully Justified Representation (FJR)}

\begin{definition}[FJR; \citealp{peters2021}]
$S$ is \emph{weakly $(\beta,T)$-cohesive} if $|S|/n\ge\cost(T)/\budget$ and $u_i(T)\ge\beta$ for
every $i\in S$. A rule satisfies \emph{FJR} if for each such $S$ some $i\in S$ has
$u_i(R(E))\ge\beta$.
\end{definition}

FJR weakens cohesion (the bundle $T$ need not be unanimously liked project-by-project, only valued
\emph{overall} by each member), so it sits strictly between EJR and the core:
\[
  \JRax \;\subset\; \PJR \;\subset\; \EJR \;\subset\; \FJR \;\subset\; \text{core}.
\]
Both PAV and Equal Shares \emph{fail} FJR (Peters et al., Examples 4--5). FJR is nonetheless
satisfiable, by the \emph{Greedy Cohesive Rule}:

\begin{definition}[Greedy Cohesive Rule (GCR)]
Start with $W=\varnothing$ and all voters active. Repeatedly find a value $\beta>0$, an active group
$S$, and $T\subseteq \projects\setminus W$ with $S$ weakly $(\beta,T)$-cohesive; among these pick
the one maximizing $\beta$ (ties to smaller $\cost(T)$). Add $T$ to $W$, deactivate $S$. Stop when
no such triple exists.
\end{definition}

\begin{theorem}[\citealp{peters2021}]
GCR satisfies FJR. Its proof uses only monotonicity of utilities, so GCR satisfies FJR even for
general monotone (non-additive) valuations.
\end{theorem}

\begin{watchout}
GCR is custom-built for FJR and pays for it: it is \emph{not} polynomial, \emph{not} priceable or
exhaustive on its own (though every GCR outcome can be \emph{completed} to a priceable one), and it
can violate \emph{laminar proportionality} that Equal Shares respects (Peters et al., Examples
6--7). For Durango, Equal Shares---EJR-satisfying, polynomial, priceable---remains the rule to
propose; FJR/GCR is the theoretical ceiling we cite, not the rule we recommend.
\end{watchout}

% ═════════════════════════════════════════════════════════════════════
\section{The DEM's distortions, in this framework}
\label{sec:distortions}

We now recast our five Durango distortions against the EJR baseline. (Full proofs live in
\texttt{notebook\_v2}; here we state each crisply in Peters et al.'s notation and flag where the
foundation sharpens it.) Throughout, $X^{\DEM}$ and $X^{\UAM}$ denote funded sets under the two
mechanisms.

\begin{proposition}[Truncation welfare loss]
\label{prop:trunc}
Because $\DEM$ forces all $\votes$ votes into one chosen district, a voter cannot express value for
projects outside it. There exist profiles where $X^{\DEM}$ is strictly Pareto-dominated in
utilitarian welfare $\sum_i u_i(\cdot)$ by an unconstrained outcome---boundary and city-wide
projects (Circleville's $A,P$) are the canonical losers.
\end{proposition}

\begin{proposition}[Noise injection via $\kappa_i$]
\label{prop:noise}
Suppose voter $i$ values only projects in $\kappa_i$ but must cast $\votes$ votes in district
$d_i$. Decompose a project's vote total as $V(c)=\text{signal}(c)+\eta(c)$, where $\eta(c)$
collects votes from voters with $c\notin\kappa_i$. Then $\E[\eta(c)]$ is independent of $c$'s
quality, and for small/low-salience projects $\eta$ can dominate signal---so vote totals are a
biased proxy for value precisely where the budget is tightest.
\end{proposition}

\begin{proposition}[Coordination game]
\label{prop:coord}
District \emph{selection} under $\DEM$ is a strategic stage: a voter's optimal district depends on
where others go (to clear within-district funding thresholds). The induced game has multiple Nash
equilibria selecting different $X^{\DEM}$; focal points (salience $s_d$) drive selection.
\end{proposition}

\begin{proposition}[Mobilization advantage]
\label{prop:mob}
An organized institution (e.g.\ a university) acts as an equilibrium selector, steering its members
into one district and displacing roughly $\votes\cdot|\text{members}|$ votes from those members'
otherwise-preferred districts---an advantage unrelated to project quality.
\end{proposition}

\begin{proposition}[Equity tradeoff: why the constraint exists]
\label{prop:equity}
Under $\UAM$, salience dominates: with $s_{d_H}>s_{d_L}$, low-salience district $d_L$ can receive
\emph{zero} funding. The $\DEM$'s per-district budgets $\budget_d$ prevent this---so the
constraint, for all its faults, is an \emph{accidental equity instrument}. Any reform must keep the
$\popshare_d$ floor it provides.
\end{proposition}

\begin{connection}
Propositions \ref{prop:trunc}--\ref{prop:equity} are exactly Peters et al.'s informal Circleville
complaints turned into propositions, plus two Durango-specific ones ($\kappa_i$ noise, mobilization)
that their approval model does not foreground. The equity tradeoff (\ref{prop:equity}) is the pivot:
it forbids the lazy reform ``just drop the constraint'' ($=\UAM$) and forces us to a rule that keeps
the floor \emph{endogenously}.
\end{connection}

% ═════════════════════════════════════════════════════════════════════
\section{How Equal Shares repairs each distortion}
\label{sec:repair}

\begin{proposition}[No truncation]
Under $\MES$ there is a single city-wide ballot: a voter assigns utility to \emph{any} project, so
boundary and city-wide projects compete on equal footing. The truncation of
Prop.~\ref{prop:trunc} cannot arise.
\end{proposition}

\begin{proposition}[No forced noise]
A voter pays only for projects with $u_i(c)>0$, i.e.\ within $\kappa_i$. No mechanism pressure
converts ignorance into votes, removing the $\eta(c)$ term of Prop.~\ref{prop:noise}.
\end{proposition}

\begin{proposition}[Reduced coordination]
There is no district to select, so the $\DEM$ selection game (Prop.~\ref{prop:coord}) disappears.
$\MES$ is still not strategyproof (below), but the multiplicity of district-selection equilibria is
gone.
\end{proposition}

\begin{proposition}[Bounded mobilization]
An institution can still concentrate utility, but $\MES$ charges each supporter from her own share
$\budget/n$: a bloc of size $k$ commands at most $k\budget/n$ of spending. Mobilization advantage is
\emph{capped by group size}, not amplified by displacement as in Prop.~\ref{prop:mob}.
\end{proposition}

\begin{theorem}[$\popshare_d$-geographic equity]
\label{thm:geoequity}
\textbf{(Coined name; result is a corollary of EJR.)} Under $\MES$, if the residents of district
$d$ can name an affordable bundle of $d$'s projects they all value, then $d$ commands at least its
population share $\popshare_d$ of the budget. Hence no district is zero-funded as it can be under
$\UAM$ (Prop.~\ref{prop:equity})---the equity the $\DEM$ provided by quota is provided by $\MES$ by
construction (Corollary~\ref{cor:district}).
\end{theorem}

\begin{naming}
We call the guarantee of Theorem~\ref{thm:geoequity} \emph{$\popshare_d$-geographic equity}. The
\emph{name} is ours; the \emph{content} is Peters et al.'s EJR specialized to the district
structure. Stating it as ``equity by construction'' is the paper's central selling point for the
reform.
\end{naming}

\begin{proposition}[Honest limitation: $\MES$ is not strategyproof]
\label{prop:notsp}
There are profiles where a voter benefits from misreporting (concrete counterexample in
\texttt{notebook\_v2}). \textbf{The operative reason is not Gibbard--Satterthwaite}---GS needs an
unrestricted preference domain, which PB does not have---but the EJR-style tension itself: a rule
delivering proportional shares cannot also be strategyproof (cf.\ Peters \& Skowron 2020 on the
limits of welfarism). We report this openly rather than overclaim.
\end{proposition}

\begin{watchout}
This is the one place where the foundation \emph{corrects} an earlier draft instinct: do not cite
Gibbard--Satterthwaite for $\MES$'s non-strategyproofness. The PB domain is restricted; the right
citation is the proportionality/welfarism literature (Peters \& Skowron 2020), not GS.
\end{watchout}

% ═════════════════════════════════════════════════════════════════════
\section{Comparison: DEM vs.\ UAM vs.\ MES}
\label{sec:comparison}

\begin{center}\small\renewcommand{\arraystretch}{1.3}
\begin{tabularx}{\linewidth}{@{}lXXX@{}}
\toprule
\textbf{Property} & \textbf{DEM} (Durango now) & \textbf{UAM} (drop constraint) & \textbf{MES} (proposed) \\
\midrule
City-wide ballot          & No (one district) & Yes & Yes \\
Boundary / city-wide goods & Orphaned & OK & OK \\
Forces uninformed votes ($\kappa_i$) & Yes ($\votes$ per district) & No & No \\
Geographic equity         & By quota $\budget_d$ & \textbf{None} (salience wins) & \textbf{By construction} ($\popshare_d$, EJR) \\
Non-geographic groups     & Silenced & Plurality only & Represented (EJR) \\
Proportionality axiom met & --- & --- (fails JR) & EJR (exact, approval) \\
Priceable                 & --- & No & Yes \\
Polynomial-time           & Yes & Yes & Yes ($O(m^2n\log n)$) \\
Strategyproof             & No & No & No (Prop.~\ref{prop:notsp}) \\
\bottomrule
\end{tabularx}
\end{center}

\begin{keyidea}
$\UAM$ fixes $\DEM$'s fragmentation but destroys equity; $\DEM$ protects equity but fragments and
adds noise. $\MES$ is the only column that is ``Yes/OK/By construction'' down the equity-and-voice
rows while staying polynomial---which is the whole argument for the reform.
\end{keyidea}

% ═════════════════════════════════════════════════════════════════════
\section{Toward the paper: the simulation \emph{(stub --- to expand)}}
\label{sec:simulation}

\begin{aside}\small
\textbf{This section is intentionally a stub in the first cut.} It records the design so the build
order is fixed; the full treatment comes once we lock the data model.
\end{aside}

Planned, mirroring Peters et al.'s experiments (their \S6) but on Durango ballots:
\begin{itemize}[leftmargin=1.4em,itemsep=2pt,topsep=2pt]
  \item \textbf{Counterfactual.} Given the same ballots, compute $X^{\DEM}$ (actual) and
  $X^{\MES}$; compare funded sets, utilitarian welfare, and the distribution of voter utility
  (Peters et al.\ find global rules give \emph{more equal} utility and usually higher total than
  separate district elections).
  \item \textbf{Utility reading.} Run both \emph{approval} and \emph{cost} utilities (\S\ref{sec:model})
  as a robustness axis.
  \item \textbf{Completion.} Compare varying-budget vs.\ perturbation vs.\ utilitarian-greedy
  completions (\S\ref{sec:axioms}).
  \item \textbf{Equity metric.} Report each district's funded share against $\popshare_d$ to show
  Theorem~\ref{thm:geoequity} bites empirically.
  \item \textbf{Synthetic ballots.} If individual ballots are unavailable, generate ballots
  consistent with published aggregates (votes per project, turnout per district).
  \item \textbf{Out of scope here:} CONAPO need-weighting and the ENIGH welfare link (follow-up).
\end{itemize}

% ═════════════════════════════════════════════════════════════════════
\appendix
\section{Notation crosswalk: \texttt{notebook\_v2} $\to$ this notebook}
\label{app:crosswalk}

\begin{center}\small\renewcommand{\arraystretch}{1.25}
\begin{tabularx}{\linewidth}{@{}lllX@{}}
\toprule
\textbf{Object} & \textbf{old (v2)} & \textbf{new} & \textbf{note} \\
\midrule
Voters                    & $N$          & $N$              & unchanged \\
Projects                  & $\mathcal{P}$ & $C$ (item $c$)   & Peters; $c$ is one project \\
Budget                    & $B$          & $b$              & Peters \\
Cost                      & $c$ (macro)  & $\cost(\cdot)$   & frees $c$ for a project \\
Per-voter votes (DEM)     & $v$          & $v$              & unchanged \\
MES threshold             & $\varrho$    & $\rho$           & dodge retired \\
EJR cohesion threshold    & ---          & $\alpha(c)$      & Peters \\
FJR weak-cohesion thresh. & ---          & $\beta$          & Peters \\
District salience         & $\alpha_d$   & $s_d$            & frees $\alpha$ \\
Information function       & $\sigma_i$   & $\kappa_i$       & frees $\sigma$ ($\sigma_c$ in EJR proof) \\
District population share & $\beta_d$    & $\pi_d$          & frees $\beta$ \\
Signal / noise            & $S(p),\eta(p)$ & $S(c),\eta(c)$ & only the project symbol changes \\
\bottomrule
\end{tabularx}
\end{center}

\section{Keystone references to add to \texttt{references.bib}}
\label{app:bib}

The load-bearing citations from Peters et al.\ for our paper are collected in
\texttt{notebooks/refs\_peters.bib} (BibTeX, ready to merge into \texttt{paper/references.bib}; the
keys match the \texttt{\textbackslash citet}/\texttt{\textbackslash citealp} calls used above). The
full list is rendered in the References below.

\nocite{stolicki2020,munagala2022,thiele1895,azizlee2020,lackner2022,sanchez2017}
\renewcommand{\bibsection}{\subsection*{References}}
\bibliographystyle{plainnat}
\bibliography{refs_peters}

\begin{aside}\small
\textbf{Status.} Tight first cut. Sections 1--7 are complete; Section~\ref{sec:simulation} is a
deliberate stub. Full proofs of Propositions~\ref{prop:trunc}--\ref{prop:notsp} remain in
\texttt{notebook\_v2} and will be migrated into this notation as we expand.
\end{aside}

\end{document}
